\documentclass[10pt,letterpaper,sans]{moderncv}
\moderncvstyle{banking}
\moderncvcolor{blue}

\usepackage[scale=0.84]{geometry}
\nopagenumbers{}
\setlength{\hintscolumnwidth}{2.35cm}

\name{Krishnaa}{Vadivel}
\title{Software Engineer | Machine Learning | Distributed Computing}
\email{srikrishnaa.careers@gmail.com}
\phone[mobile]{516-853-5663}
\social[linkedin]{sri-krishnaa-ece-phd}
\extrainfo{\href{https://krishnaa423.github.io/personal-website/}{Website} \textbar{} \href{https://github.com/krishnaa423}{GitHub: krishnaa423} \textbar{} \href{https://leetcode.com/u/krishnaa42342}{LeetCode: krishnaa42342} \textbar{} \href{https://hub.docker.com/u/krishnaa42342}{Docker Hub: krishnaa42342}}

\begin{document}
\makecvtitle

\section{Profile}
\cvitem{}{Software engineer and PhD researcher focused on machine learning, distributed computing, and numerical software. Strong background in Python and C++ development for simulation, automation, GPU acceleration, data systems, and infrastructure spanning research and production-adjacent environments.}

\section{Research Experience}
\cventry{2021--present}{Research Assistant / PhD Researcher}{Yale University, Electrical Engineering}{}{}{\begin{itemize}%
\item Developed graph neural networks and automatic-differentiation models in Python to predict material properties including optical emission spectra.
\item Developed scientific software for first-principles and many-body simulation, combining Python and C++ with MPI, OpenMP, PETSc, SLEPc, CUDA, ROCm, and CMake.
\item Built distributed and accelerator-enabled applications for HPC environments including Frontier, Frontera, and Perlmutter, with experience spanning NVIDIA and AMD GPU systems on national-lab clusters.
\item Worked across debugging, compiler-aware builds, documentation, and reproducible computational software for research-scale simulations.
\item Extended into ML-assisted and retrieval-driven automation with LangChain, OpenAI API models, and local Ollama models for coding and research setup.
\item Co-authored a 2025 \textit{Journal of the American Chemical Society} paper and presented APS talks in 2024, 2025, and 2026 on self-trapped exciton formation and related first-principles work.
\end{itemize}}

\section{Professional Experience}
\cventry{2019--2021}{Software and Embedded Systems Engineer}{Probe Technologies Inc. / Weatherford Wireline Services}{}{}{\begin{itemize}%
\item Built parallel applications with CUDA and OpenCL fallback for large acoustic waveform and sensor datasets used in engineering diagnostics and field-tool analysis.
\item Developed full-stack server applications with ASP.NET and Microsoft SQL Server for real-time control of company tool testing, data collection, and data processing.
\item Developed C\# GUI applications for preparing customer inputs and interacting with deployed hardware systems.
\item Worked close to embedded and instrumentation systems, giving practical experience with software operating alongside deployed hardware and field constraints.
\end{itemize}}
\cventry{2018--2019}{Photonics Technical Writer}{FindLight}{}{}{\begin{itemize}%
\item Produced technical content for optics and photonics topics, demonstrating strong technical communication across specialized engineering subject matter.
\end{itemize}}

\section{Selected Projects}
\cventry{}{Machine Learning Models for Material Properties}{}{}{}{\begin{itemize}%
\item Developed PyTorch Geometric graph neural networks and automatic-differentiation models for molecular-property prediction and optical-properties calculations in material systems.
\end{itemize}}
\cventry{}{Agentic Coding and Scientific Workflow Generation}{}{}{}{\begin{itemize}%
\item Developed LangChain-based retrieval-driven tooling using OpenAI API models and local Ollama models to read documentation and academic papers, then generate structured starter inputs and workflow artifacts.
\end{itemize}}
\cventry{}{Container Orchestration for First-Principles Computing}{}{}{}{\begin{itemize}%
\item Developed CUDA-ready and cross-environment containers for first-principles software stacks and published them to Docker Hub for portable scientific computing.
\end{itemize}}
\cventry{}{Matrix-Multiply AI Accelerator}{}{}{}{\begin{itemize}%
\item Developed a Verilog-based matrix-multiply accelerator project for AI-oriented compute pipelines and hardware-software performance exploration.
\end{itemize}}

\section{Education}
\cventry{2021--present}{PhD, Electrical Engineering}{Yale University}{}{}{Ongoing work combining scientific computing, high-performance workflows, and computational materials research.}
\cventry{2023}{MPhil, Electrical Engineering}{Yale University}{}{}{}
\cventry{2023}{MS, Electrical Engineering}{Yale University}{}{}{}
\cventry{2017--2018}{MEng, Electrical and Computer Engineering}{Cornell University}{}{}{}
\cventry{2013--2017}{BS, Electrical and Computer Engineering}{Cornell University}{}{}{}

\section{Selected Coursework}
\cvitem{Computing Foundations}{Theoretical Computer Science; Probabilistic Machine Learning}
\cvitem{Scientific Foundations}{Graduate Quantum Mechanics; Solid State Physics; Many-Body Quantum Physics}
\cvitem{Engineering Foundations}{Signals and systems, embedded systems, semiconductor physics, and numerical/scientific computing practice}

\section{Technical Skills}
\cvitem{Languages}{Python, C, C++, JavaScript, Bash, PowerShell, SQL}
\cvitem{Scientific Computing}{MPI, OpenMP, PETSc, SLEPc, mpi4py, petsc4py, slepc4py, NumPy, pandas, SciPy, SymPy, matplotlib}
\cvitem{Systems}{Linux command line, CMake, Docker, Docker Compose, Kubernetes, gcloud}
\cvitem{Platforms}{CUDA, ROCm, Frontier, Frontera, Perlmutter}
\cvitem{ML / Automation}{PyTorch, PyTorch Geometric, scikit-learn, LangChain, OpenAI API, Ollama, graph neural networks, transformers, agentic automation}

\end{document}
